% Abstract: 100-150 Wörter (Folie 16). Problem - Vorgehen - Kernergebnis.
\begin{abstract}
Vortrainierte Vision Transformer (ViTs) lassen sich gut auf neue Aufgaben übertragen, doch alle
Parameter anzupassen ist teuer. Wir untersuchen, wie effizient sich ein vortrainierter ViT-B/16 an die
Klassifikation der 25 Hunderassen von Oxford-IIIT Pet anpassen lässt. Mit einheitlichem Protokoll
und drei Seeds vergleichen wir Linear Probing, Low-Rank Adaptation (LoRA) mit verschiedenen Rängen und
Ziel-Layern, volles Fine-Tuning und zwei CNNs. LoRA mit Rang 8 erreicht die
höchste Test-Accuracy (96,2\,\%; die auf Validation gewählte Variante $r=4$ 96,1\,\%) bei 0,36\,\% trainierbaren Parametern, doch Linear Probing mit nur
0,02\,\% ist innerhalb der Seed-Streuung gleich gut. Volles Fine-Tuning liegt 0,8 Prozentpunkte darunter,
das beste CNN 1,4 Punkte unter dem besten ViT. Zusatzexperimente liefern Hinweise auf die hohe Genauigkeit: 22 der 25
Rassen sind ImageNet-Klassen, und die ursprünglichen ImageNet-Köpfe erkennen sie ganz ohne Training zu
96--97\,\% (Auswahl unter diesen 22). Bei Daten so nah am Pretraining ist ein neuer linearer Kopf die
effizienteste Anpassung.
\end{abstract}
